\begin{figure}[H]
\centering
\resizebox{0.88\linewidth}{!}{%
\begin{tikzpicture}[>=Stealth,
  stream/.style={draw=black!55, line width=0.8pt},
  streamline/.style={draw=black!35, densely dashed, line width=0.55pt},
  section/.style={draw=black!65, line width=0.75pt},
  velocity/.style={->, draw=blue!70!black, line width=1.05pt, >={Stealth[length=2.8mm]}},
  pressure/.style={draw=red!70!black, line width=1.05pt},
  annot/.style={font=\footnotesize\sffamily, fill=white, inner sep=1.2pt},
]
  % 流管：速度沿程增加，因此截面积沿程减小
  \draw[stream] (-0.2,1.35) .. controls (1.7,1.28) and (2.6,1.02) .. (3.2,0.9)
                .. controls (4.4,0.67) and (5.6,0.57) .. (6.8,0.55);
  \draw[stream] (-0.2,-1.35) .. controls (1.7,-1.28) and (2.6,-1.02) .. (3.2,-0.9)
                .. controls (4.4,-0.67) and (5.6,-0.57) .. (6.8,-0.55);
  \draw[streamline] (-0.2,0.68) .. controls (1.7,0.64) and (2.6,0.51) .. (3.2,0.45)
                     .. controls (4.4,0.34) and (5.6,0.29) .. (6.8,0.28);
  \draw[streamline] (-0.2,-0.68) .. controls (1.7,-0.64) and (2.6,-0.51) .. (3.2,-0.45)
                     .. controls (4.4,-0.34) and (5.6,-0.29) .. (6.8,-0.28);

  % 三个截面
  \draw[section] (0,1.35) -- (0,-1.35);
  \draw[section, line width=1.5pt, fill=black!10, fill opacity=.45] (3.2,0.9) -- (3.2,-0.9);
  \draw[section] (6.6,0.56) -- (6.6,-0.56);
  \node[annot, above=2pt, align=center] at (0,1.35) {截面0\\$A_0$};
  \node[annot, above=2pt, align=center] at (3.2,0.9) {桨盘截面1\\$A$};
  \node[annot, above=2pt, align=center] at (6.6,0.56) {截面2\\$A_2$};

  % 速度箭头，长度随速度增加
  \draw[velocity] (-0.85,0) -- (0.45,0) node[annot, midway, above=2pt] {$v_0$};
  \draw[velocity] (2.15,0) -- (3.75,0) node[annot, midway, above=2pt] {$v_p=v_0+v_i$};
  \draw[velocity] (4.65,0) -- (6.75,0) node[annot, midway, above=2pt] {$v_2=v_0+2v_i$};
  \node[font=\scriptsize\sffamily, text=black!60, fill=white, inner sep=1pt] at (3.2,-1.16) {无限薄桨盘};

  % 独立压强示意，避免与流管重叠
  \draw[->, draw=black!55, line width=.55pt] (-0.1,-2.35) -- (7.05,-2.35) node[right, font=\scriptsize\sffamily] {$x$};
  \draw[->, draw=black!55, line width=.55pt] (-0.1,-2.35) -- (-0.1,-1.30)
    node[above, font=\scriptsize\sffamily] {$p-p_\infty$};
  \draw[pressure, dashed] (-0.1,-1.88) -- (7.0,-1.88);
  \node[font=\scriptsize\sffamily, text=red!70!black, fill=white, inner sep=1pt] at (0.45,-1.73) {$p_\infty$};
  \draw[pressure] (0,-1.88) .. controls (1.5,-1.9) and (2.55,-2.12) .. (3.18,-2.16);
  \draw[pressure] (3.2,-1.42) .. controls (4.0,-1.48) and (5.35,-1.75) .. (6.6,-1.88);
  \draw[pressure, ->] (3.2,-2.16) -- (3.2,-1.42);
  \node[annot, text=red!70!black, anchor=west] at (3.38,-1.7) {$\Delta p=T/A$};
\end{tikzpicture}%
}
\caption{理想桨盘的流管与压强分布示意。由于$v_0<v_p<v_2$，不可压定常流的流管沿流向收缩（$A_0>A>A_2$）；速度穿过无限薄桨盘时连续，而静压在盘面发生跃变$\Delta p=T/A$，并在远尾流中恢复至环境压强。图中 $x$ 为桨盘局部流向轴，不代表构型主图的相机左右；几何尺寸与箭头长度仅作定性示意。}
\label{fig:actuator_disk}
\end{figure}
